\documentclass[11pt,a4paper]{article}

% ========== PACKAGES ==========
\usepackage{graphicx}
\usepackage{booktabs}
\usepackage{hyperref}
\usepackage{amsmath}
\usepackage{amssymb}
\usepackage{microtype}
\microtypesetup{expansion=false}
\usepackage{xcolor}
\usepackage{enumitem}
\usepackage{caption}

% ========== BASIC TWO-COLUMN SETUP ==========
\usepackage[margin=1in]{geometry}
\usepackage{multicol}
\usepackage{hyphenat}
\setlength{\columnsep}{0.6cm}

% ========== TITLE & AUTHORS ==========
\title{\vspace{-1.5cm}Why RAG Fails: A Taxonomy of Failure Modes in Retrieval-Augmented Generation}

\author{Your Name \\
Your University / Affiliation \\
\texttt{your.email@example.com}
}

\date{}

\begin{document}

\maketitle

\begin{multicols}{2}

% ========== ABSTRACT ==========
\begin{abstract}
Retrieval-Augmented Generation (RAG) has become a standard paradigm for grounding large language models in external knowledge, yet RAG systems still fail in production in ways that are poorly understood and rarely standardized. Existing work focuses on architectures and accuracy benchmarks, leaving practitioners without a common language for describing and diagnosing RAG failures. In this paper, we propose a failure-centric taxonomy of five modes in RAG systems: Retrieval Blindness, Retrieval Noise, Evidence Ignorance, Evidence Over-Reliance, and Latency-Induced Failure. For each mode, we define typical patterns and causes and illustrate them with concrete scenarios. We further introduce a practical diagnosis guide that maps observable symptoms to failure modes and candidate mitigations, and illustrate the taxonomy through four enterprise case study scenarios. Building on this conceptual framework, we contribute a \textbf{preliminary empirical validation} of three of the five modes---Retrieval Blindness, Retrieval Noise, and Evidence Ignorance---using two fixed generator models (Claude Sonnet 4.6 and Llama 3.3 70B Instruct) under an identical retrieval pipeline, prompt template, and evaluation protocol, with controlled clean and stress conditions drawn from the RGB and Lost-in-the-Middle public benchmarks. We report token-level F1, retrieval recall/precision/nDCG, exact match, cost, and latency across 900 real API calls (450 per model), together with a manual error-categorization pass on a stratified sample of 70 outputs. Both models show substantial, consistent degradation from clean to stress conditions across all three tested modes, supporting the taxonomy's claims for \textbf{these three modes only}---Retrieval Blindness, Retrieval Noise, and Evidence Ignorance---for these two models on this dataset; the remaining two modes (Evidence Over-Reliance, Latency-Induced Failure), a third model (GPT-4.1), and broader taxonomy-wide validation all remain future work. We are explicit that this is a small-scale, two-model, single-dataset-family preliminary study, not a full validation of the taxonomy, and we detail its limitations alongside the results. Our aim is to establish both a shared vocabulary/diagnostic tool for RAG failures and a reproducible empirical basis that future work can extend.
\end{abstract}

% ========== SECTION 1: INTRODUCTION ==========
\section{Introduction}

Retrieval-Augmented Generation (RAG) has rapidly emerged as a dominant paradigm for addressing the knowledge limitations of large language models (LLMs). By integrating external retrieval mechanisms, systems such as Bing Chat, Perplexity, and numerous enterprise pipelines now ground their outputs in up-to-date, verifiable information \cite{lewis2020rag}. The promise is compelling: combine the fluency of LLMs with the factual accuracy of search.

However, the gap between promise and reality is substantial. RAG systems fail in production with alarming regularity. They retrieve correct documents yet hallucinate contradictory answers. They ignore retrieved context entirely. They amplify noise and present it as fact. Despite growing awareness of these issues, the community lacks a systematic vocabulary to describe \textit{how} and \textit{why} RAG fails.

Existing taxonomies in the RAG literature focus overwhelmingly on \textit{architectures}---classifying systems as naive RAG, advanced RAG, or modular RAG \cite{gao2023ragsurvey}. While useful for system design, these categorizations offer no guidance when a deployed system produces incorrect output. A practitioner observing a failure cannot consult the architectural literature to diagnose the root cause.

In this paper, we address this gap by proposing a failure-centric taxonomy for RAG, and by taking a first empirical step toward validating it. Our goal is to provide a conceptual framework, a practical diagnosis guide, and a reproducible empirical basis that can be adopted and extended by researchers and practitioners as a standard way to describe, debug, and test for RAG failures. Our contributions are:
\begin{enumerate}[leftmargin=*,nosep]
\item A taxonomy of five failure modes in RAG systems: Retrieval Blindness, Retrieval Noise, Evidence Ignorance, Evidence Over-Reliance, and Latency-Induced Failure.
\item Concrete examples and causes for each mode, synthesizing observations from RAG evaluation and LLM reliability studies.
\item A diagnosis guide that maps observable symptoms to failure modes and suggests mitigations, intended as a reusable tool for RAG practitioners.
\item An illustrative enterprise case study with four scenarios showing how the taxonomy and diagnosis guide apply in practice.
\item \textbf{A preliminary empirical validation} of three of the five modes (Retrieval Blindness, Retrieval Noise, and Evidence Ignorance) using two fixed generator models under an identical retrieval pipeline and evaluation protocol, with controlled clean and stress conditions, real API-measured metrics, and a manual error analysis on a stratified sample of outputs.
\item A discussion of failure-aware evaluation protocols that complement existing accuracy-centric RAG benchmarks, informed by what we learned running the preliminary study above.
\end{enumerate}

This paper originated as a conceptual and practitioner-oriented position paper focused on defining failure modes and diagnostic workflows. This version (v2) extends that foundation with a preliminary, small-scale empirical study of three of the five failure modes. We are explicit throughout about which claims remain conceptual (Evidence Over-Reliance, Latency-Induced Failure, and the taxonomy's completeness) and which are now empirically supported for the specific models, dataset, and conditions we tested.

The remainder of this paper is organized as follows.
 Section~\ref{sec:related} reviews related work. Section~\ref{sec:taxonomy} presents the failure mode taxonomy. Section~\ref{sec:diagnosis} provides the diagnosis guide and usage workflow. Section~\ref{sec:empirical} presents our preliminary empirical validation. Section~\ref{sec:casestudy} presents an enterprise RAG case study. Section~\ref{sec:evaluation} discusses implications for evaluation, and Section~\ref{sec:conclusion} concludes.

% ========== SECTION 2: RELATED WORK ==========
\section{Related Work}
\label{sec:related}

\subsection{RAG Architectures}

Since the seminal work of Lewis et al.~\cite{lewis2020rag}, RAG has evolved through several architectural paradigms. Gao et al.~\cite{gao2023ragsurvey} provide a comprehensive survey, categorizing systems into naive RAG, advanced RAG, and modular RAG. While valuable for system builders, these taxonomies classify \textit{design choices}, not \textit{failure behaviors}.

\subsection{RAG Evaluation}

Several benchmarks measure RAG performance. RAGAS~\cite{ragas2024} and RGB~\cite{rgb2024} evaluate faithfulness, relevance, and answer correctness. However, these benchmarks primarily assess performance under clean retrieval conditions and do not systematically probe resilience to specific failure modes or stress test robustness under noisy, incomplete, or adversarial retrieval.

\subsection{Reliability and Robustness}

Broader work on LLM reliability has identified hallucination~\cite{ji2023hallucination} as a persistent risk in generative systems. Within RAG, studies show models sometimes ignore provided context~\cite{liu2024lost} or are vulnerable to adversarial documents and outdated knowledge bases. Our taxonomy synthesizes these isolated findings into a unified framework centered on failure modes rather than architectures or metrics.

\subsection{Failure-Oriented RAG Studies}

Barnett et al.~\cite{barnett2024seven} present an experience report identifying seven failure points from RAG deployments in research, education, and biomedical domains. Their work provides important engineering lessons and emphasizes that RAG robustness must be validated during operation. Mindful-RAG~\cite{agrawal2024mindful} studies eight failure points in knowledge-graph-based RAG, with particular attention to query intent and context gathering. These studies are the closest prior work to our paper. Our contribution differs in offering a compact five-mode taxonomy centered on retrieval--generation misalignment, together with a lightweight diagnosis workflow that maps observable symptoms to candidate mitigations. We also include Latency-Induced Failure as a first-class production failure mode and illustrate the taxonomy through enterprise RAG scenarios.

\textbf{Gap.} Existing work identifies important RAG failure points, but there remains a need for a compact taxonomy focused on retrieval--generation misalignment and a lightweight diagnosis workflow that practitioners can use to identify and mitigate failures in deployed systems.

% ========== SECTION 3: THE FAILURE MODE TAXONOMY ==========
\section{The Failure Mode Taxonomy}
\label{sec:taxonomy}

We propose a taxonomy of five failure modes in RAG systems. These modes are not mutually exclusive; real-world failures often involve multiple categories. Each mode captures a distinct way in which retrieval and generation can misalign.

\subsection{Retrieval Blindness}

\textbf{Definition.} Retrieval Blindness occurs when the retrieval component fails to surface any useful evidence for the query. The answer is effectively generated ``in the dark.''

\textbf{Typical pattern.} The corpus does not contain the needed information, or the retriever cannot bridge the gap between user language and document language.

\textbf{Example.} A legal RAG assistant is asked about a newly introduced regulation that is not yet indexed; the system returns generic legal advice with no corpus-specific citations, despite multiple retrieval attempts.

\subsection{Retrieval Noise}

\textbf{Definition.} Retrieval Noise occurs when the retriever returns content that is irrelevant, misleading, or off-topic, contaminating the context fed to the LLM.

\textbf{Typical pattern.} Retrieved passages share surface-level tokens with the query but are semantically mismatched, leading to answers that mix correct and spurious details.

\textbf{Example.} A medical RAG system asked about ``ACE inhibitors and kidney function'' retrieves advertising copy and unrelated patient stories due to keyword overlap, causing the model to blend factual pharmacology with anecdotal noise.

\subsection{Evidence Ignorance}

\textbf{Definition.} Evidence Ignorance occurs when retrieval succeeds, but the LLM ignores the provided context and defaults to parametric knowledge or hallucinations.

\textbf{Typical pattern.} Removing the retrieved context does not significantly change the answer; the model relies on its internal priors instead of grounded evidence, often due to long-context effects or weak instructions.

\textbf{Example.} In a long-context setting, relevant passages about a company policy appear in the middle of the context window, yet the model answers from outdated parametric knowledge, illustrating lost-in-the-middle behavior~\cite{liu2024lost}.

\subsection{Evidence Over-Reliance}

\textbf{Definition.} Evidence Over-Reliance occurs when the LLM trusts retrieved content too readily, propagating incorrect, outdated, or adversarial information.

\textbf{Typical pattern.} Updating or correcting the underlying corpus changes the failure, indicating that the model faithfully followed bad evidence rather than hallucinating independently.

\textbf{Example.} A finance RAG assistant retrieves an outdated prospectus and repeats superseded interest rates as current, without cross-checking against newer documents or external signals.

\subsection{Latency-Induced Failure}

\textbf{Definition.} Latency-Induced Failure occurs when retrieval cannot complete within time or resource constraints, leading the system to answer with partial, stale, or missing context.

\textbf{Typical pattern.} Under high load, multi-hop or cross-source retrieval is truncated; the model generates short, incomplete answers or silently falls back to unguided parametric responses.

\textbf{Example.} An enterprise RAG dashboard under peak traffic times out on secondary retrieval steps, returning summaries that omit key risk factors because the relevant documents were never retrieved in time.

\paragraph{Multi-mode failures.}
In practice, many real-world failures involve combinations of these modes rather than a single category. For example, Retrieval Noise and Evidence Over-Reliance often co-occur when a retriever surfaces outdated or off-topic documents and the model trusts them without filtering, as in Scenario 2 of our case study. We recommend labeling failures with all applicable modes and then prioritizing mitigations that address upstream causes (e.g., improving retrieval quality and corpus hygiene) before downstream behaviors (e.g., prompt design or generation strategies). This multi-label perspective reflects the fact that RAG failures arise from interactions between components, and helps practitioners reason about compound failure patterns in a structured way.

Latency-induced failures are more infrastructural than the other four modes, but they directly affect whether retrieval and generation can align at all. When retrieval is truncated or skipped due to latency constraints, the system is effectively forced into Retrieval Blindness or Evidence Ignorance, answering from partial or missing evidence. In production RAG deployments, latency budgets and system load are common root causes of misalignment between retrieved context and generated answers, so we include Latency-Induced Failure as a first-class mode in our taxonomy.

% ========== SECTION 4: DIAGNOSIS GUIDE ==========
\section{Diagnosis Guide}
\label{sec:diagnosis}

To make the taxonomy actionable, we provide a diagnosis guide that maps observable symptoms in RAG systems to likely failure modes and candidate mitigations. Practitioners can use the following summary as a checklist when debugging RAG systems.

\subsection*{Summary of failure modes}

\textbf{Retrieval Blindness.}
Observable symptom: Answers are generic and vague, with few specific facts tied to the corpus; adding more documents does not help.
Likely cause: Retriever fails to surface useful evidence, or the corpus lacks the needed information.
Candidate mitigations: Expand or update the knowledge base; improve chunking granularity; adapt or fine-tune embeddings to the domain; add query rewriting or template-based queries.

\textbf{Retrieval Noise.}
Observable symptom: Answers mix correct and irrelevant information; retrieved context appears off-topic or misleading, often with large top-$k$ and low precision.
Likely cause: Semantic similarity failures, keyword collisions, mixed-topic chunks.
Candidate mitigations: Reduce top-$k$; introduce re-ranking using cross-encoders; enforce tighter filtering on metadata; fix chunk boundaries to avoid mixed-topic segments.

\textbf{Evidence Ignorance.}
Observable symptom: Correct documents are retrieved, but the model's answer contradicts or ignores them; removing context barely changes the answer.
Likely cause: Parametric memory dominance, long-context effects, weak instructions to use evidence.
Candidate mitigations: Strengthen instructions to cite and use retrieved evidence; place key passages near the beginning of the context; reduce context length to avoid lost-in-the-middle effects; tune prompts to discourage parametric guesses.

\textbf{Evidence Over-Reliance.}
Observable symptom: Answers confidently repeat incorrect or outdated facts present in the retrieved documents; updating the corpus changes the failure.
Likely cause: Lack of verification or cross-checking; outdated or adversarial content in the corpus.
Candidate mitigations: Add verification or cross-checking modules; mark and filter outdated content; incorporate multiple sources and consistency checks; use uncertainty estimation or citations with confidence scores.

\textbf{Latency-Induced Failure.}
Observable symptom: Under load or strict latency budgets, answers become short, incomplete, or ignore context when timeouts occur.
Likely cause: Retrieval and orchestration exceeding latency budgets; multi-hop retrieval truncated; inadequate fallback behavior.
Candidate mitigations: Profile and optimize retrieval latency; cache frequent queries; design graceful degradation strategies (e.g., fallback summaries or lower-cost retrieval); adjust timeouts and budget allocation across retrieval and generation.

Practitioners can treat this summary as a checklist: observe a failure, map it to one or more modes, and then apply the corresponding mitigations when iterating on RAG system design.

\subsection{Using the Taxonomy in Practice}

Practitioners can apply our taxonomy and diagnosis guide in a lightweight debugging workflow:
\begin{enumerate}[leftmargin=*,nosep]
\item \textbf{Log failures with symptoms.} When a RAG system produces an incorrect or unsatisfactory answer, record the query, retrieved documents, answer, and a short description of the observable symptom (e.g., ``generic answer'', ``mix of relevant and irrelevant facts'').
\item \textbf{Map failures to modes.} Use the summary above to map each failure to one or more modes based on its symptom pattern. This creates a standardized label for discussing the failure across teams.
\item \textbf{Track mode frequencies.} Aggregate failures over time by mode (e.g., percentage of failures due to Retrieval Noise vs. Evidence Ignorance). This highlights which parts of the RAG pipeline contribute most to observed issues.
\item \textbf{Prioritize mitigations.} For the most frequent modes, apply the candidate mitigations listed in the summary---for example, investing in better retrieval for Retrieval Noise or stronger prompting for Evidence Ignorance.
\end{enumerate}
This process can be integrated into existing RAG debugging and monitoring workflows with minimal overhead, and provides a shared vocabulary for postmortems, design discussions, and robustness reports.


% ========== SECTION: PRELIMINARY EMPIRICAL VALIDATION ==========
\section{Preliminary Empirical Validation}
\label{sec:empirical}

Sections~\ref{sec:taxonomy} and~\ref{sec:diagnosis} presented Retrieval Blindness, Retrieval Noise, Evidence Ignorance, Evidence Over-Reliance, and Latency-Induced Failure as conceptual categories, illustrated with hypothetical examples. This section reports a preliminary empirical study of three of the five modes---Retrieval Blindness, Retrieval Noise, and Evidence Ignorance---using controlled conditions, two fixed generator models, and real, cost- and latency-logged API calls. Evidence Over-Reliance and Latency-Induced Failure remain conceptual in this version of the paper; we did not design or run controlled experiments for either mode here.

\textbf{Scope and honesty about limits.} This is a \textit{preliminary} study: two generator models, one dataset family per failure mode, and 30 base questions per mode. It demonstrates that the three tested failure modes are real, measurable, and reproducible for these two models under these conditions---it does not establish that the full five-mode taxonomy is validated, nor that these findings generalize to other models, retrievers, or domains. Section~\ref{sec:empirical-limitations} details these limitations explicitly.

\subsection{Experimental Design}
\label{sec:empirical-design}

\textbf{Dataset.} We use two existing public RAG evaluation datasets rather than hand-written questions. For Retrieval Blindness and Retrieval Noise, we draw 30 questions per mode from the RGB benchmark~\cite{rgb2024} (English subset, \textit{en\_refine.json}, retrieved from the authors' repository, commit \textit{65ec39e}), which pairs short-answer factual questions with curated positive and negative (noise) passages. For Evidence Ignorance, we draw 30 questions from the Lost-in-the-Middle Natural Questions--Open resource~\cite{liu2024lost} (\textit{nq-open-10\_total\_documents\_gold\_at\_4.jsonl.gz}, commit \textit{29b8a6d}), which places a single gold passage at a controlled position among 10 total distractor documents. All 90 base questions (30 per mode) are fixed once and reused identically across both models; no question was selected, discarded, or reweighted after seeing any model's output.

\textbf{Retriever and embeddings.} We use a single deterministic flat cosine-similarity retriever (no approximate search, no managed vector database) over \textit{sentence-transformers/all-MiniLM-L6-v2}~\cite{reimers2019minilm} embeddings (CPU inference, L2-normalized), with top-$k=5$ for every query in every condition. RGB and Lost-in-the-Middle ship pre-chunked passages; we apply no additional chunking. The retriever, embedding model, chunking, top-$k$, and random seed (42) are identical across both generator models and all conditions---only the generator model differs between runs.

\textbf{Prompt template and generation settings.} Every call uses the same prompt template (context block followed by the question, instructing the model to answer only from the provided context and to say ``I don't know based on the available information'' rather than guess), the same generation temperature (0.0), and the same maximum output length (200 tokens), for both models and all conditions.

\textbf{Conditions.} For each failure mode we construct a clean condition and one or more stress conditions, holding the underlying question and gold answer fixed:
\begin{itemize}[leftmargin=*,nosep]
\item \textit{Retrieval Blindness}: \textbf{clean} (gold passage retrievable) vs.\ \textbf{blind} (gold passage withheld from the retrieval index for that query, forcing the model to answer from whatever else is retrieved).
\item \textit{Retrieval Noise}: \textbf{clean} (top-5 retrieval from the curated positive passages) vs.\ \textbf{noisy} (RGB's curated lexically/semantically similar distractor passages injected into the candidate pool, so the gold passage must compete with real, dataset-provided noise for the top-5 slots).
\item \textit{Evidence Ignorance}: \textbf{top} (gold passage at position 1 of 10), \textbf{middle} (gold passage at position 4 of 10, the Lost-in-the-Middle configuration), and \textbf{removed} (gold passage dropped, only the 9 distractors remain).
\end{itemize}
Clean and \textit{top} conditions use a single trial per question; the four stress conditions (\textbf{blind}, \textbf{noisy}, \textbf{middle}, \textbf{removed}) use three trials per question at temperature 0.0 to check residual run-to-run variation even under nominally deterministic decoding. This yields $30 + 90 = 120$ generation calls per failure mode per model (450 total per model across all three modes).

\textbf{Models.} We use exactly two of the three originally planned fixed generator models: \textit{anthropic/claude-sonnet-4.6}~\cite{anthropic2026sonnet46} and \textit{meta-llama/llama-3.3-70b-instruct}~\cite{meta2024llama33}, both accessed through the OpenRouter API~\cite{openrouter2026} (Claude: routed at OpenRouter's standard \$3/M input, \$15/M output token pricing with no per-token markup over the native Anthropic rate; Llama: OpenRouter's \$0.10/M input, \$0.32/M output pricing). All 900 generation calls (450 per model) were made on 2026-08-15 through 2026-08-17. GPT-4.1~\cite{openai2025gpt41} was in scope for the original design but is deferred to future work; Section~\ref{sec:empirical-limitations} and Section~\ref{sec:conclusion} explain why.

\textbf{Metrics.} We report token-level F1 (word-overlap F1 between generated and reference answer) and exact match as answer-quality metrics; retrieval recall@5, precision@5, and nDCG@5 (computed against each dataset's own gold-passage labels) as retrieval-quality metrics; and total input/output tokens, per-call latency, and total API cost as efficiency metrics. We flag exact match as a known limitation for this dataset combination: many gold answers include full proper names or exact phrasings (e.g., ``Hero Beauregard Fiennes-Tiffin'') that a substantively correct answer (``Hero Fiennes-Tiffin'') will not match under strict string equality, so exact match systematically understates correctness here; token F1 is the more informative of the two answer-quality metrics for this study, and we report exact match for completeness rather than as a primary measure. A second, distinct case in the same LITM subset compounds this: item \texttt{litm-2945d0775c04} asks a two-part question (``when did the Soviet Union first gain control of parts of Poland and the Baltic republics'') but its gold answer, ``14 June 1940,'' covers only the Baltic half (Poland was 17 September 1939); both models correctly identified both dates in every condition yet scored exact match 0 by construction, independent of correctness, which we verified against external historical sources.

\subsection{Results}
\label{sec:empirical-results}

Table~\ref{tab:claude-results} and Table~\ref{tab:llama-results} report per-condition metrics for Claude Sonnet 4.6 and Llama 3.3 70B Instruct respectively, averaged over all trials. $n$ is the number of generation calls in that row (30 for single-trial conditions, 90 for three-trial stress conditions).

\begin{center}
\scriptsize
\setlength{\tabcolsep}{2.5pt}
\begin{tabular}{@{}llrrrrr@{}}
\toprule
Mode & Cond. & $n$ & F1 & Rec. & P@5 & nDCG \\
\midrule
RB & clean & 30 & 0.108 & 0.767 & 0.153 & 0.544 \\
RB & blind & 90 & 0.028 & 0.000 & 0.000 & 0.000 \\
RN & clean & 30 & 0.117 & 1.000 & 1.000 & 1.000 \\
RN & noisy & 90 & 0.066 & 0.633 & 0.127 & 0.479 \\
EI & top & 30 & 0.067 & 1.000 & 0.200 & 1.000 \\
EI & middle & 90 & 0.067 & 1.000 & 0.200 & 0.500 \\
EI & removed & 90 & 0.032 & 0.000 & 0.000 & 0.000 \\
\bottomrule
\end{tabular}
\captionof{table}{Claude Sonnet 4.6, 450/450 real calls, 0 errors. RB = Retrieval Blindness, RN = Retrieval Noise, EI = Evidence Ignorance. Exact match was 0.000 in every row and is omitted from this table (see text). Total: 400{,}538 input / 41{,}438 output tokens, avg.\ latency 2{,}528\,ms, cost \$1.82.}
\label{tab:claude-results}
\end{center}

\begin{center}
\scriptsize
\setlength{\tabcolsep}{2.5pt}
\begin{tabular}{@{}llrrrrrr@{}}
\toprule
Mode & Cond. & $n$ & F1 & EM & Rec. & P@5 & nDCG \\
\midrule
RB & clean & 30 & 0.260 & 0.067 & 0.767 & 0.153 & 0.544 \\
RB & blind & 90 & 0.029 & 0.000 & 0.000 & 0.000 & 0.000 \\
RN & clean & 30 & 0.446 & 0.233 & 1.000 & 1.000 & 1.000 \\
RN & noisy & 90 & 0.230 & 0.056 & 0.633 & 0.127 & 0.479 \\
EI & top & 30 & 0.189 & 0.000 & 1.000 & 0.200 & 1.000 \\
EI & middle & 90 & 0.182 & 0.000 & 1.000 & 0.200 & 0.500 \\
EI & removed & 90 & 0.063 & 0.000 & 0.000 & 0.000 & 0.000 \\
\bottomrule
\end{tabular}
\captionof{table}{Llama 3.3 70B Instruct, 450/450 real calls, 0 errors. Unlike Claude, Llama registered nonzero exact match (EM) on several conditions, so EM is shown here. Total: 367{,}827 input / 16{,}964 output tokens, avg.\ latency 2{,}674\,ms, cost \$0.042.}
\label{tab:llama-results}
\end{center}

\textbf{Retrieval recall, precision@5, and nDCG@5 are identical between the two models within each condition.} This is expected and not a coincidence: both models share the exact same retriever, embedding model, and top-$k$, so retrieval quality is a property of the pipeline and the condition, not of the generator. The retrieval-side numbers isolate how much relevant evidence each condition makes available; the F1/EM columns then show how each model uses (or fails to use) that evidence.

\textbf{Retrieval Blindness.} Withholding the gold passage collapses retrieval recall/precision/nDCG to exactly 0 for both models (by construction) and drives token F1 down by roughly 4$\times$ for Claude (0.108 $\rightarrow$ 0.028) and 9$\times$ for Llama (0.260 $\rightarrow$ 0.029). Both models retain some non-zero F1 in the blind condition, indicating partial credit from generic or partially-relevant guesses rather than complete failure; our error analysis (Section~\ref{sec:empirical-erroranalysis}) shows this residual score is a mix of well-calibrated refusals, confidently wrong grounded-in-distractor guesses, and occasional correct answers drawn from parametric knowledge rather than the (withheld) retrieved evidence.

\textbf{Retrieval Noise.} Injecting RGB's curated distractor passages drops recall from 1.000 to 0.633 and precision@5 from 1.000 to 0.127, confirming that noise passages measurably displace gold passages from the top-5 window rather than merely sitting alongside them unused. Token F1 falls by about 44\% for Claude (0.117 $\rightarrow$ 0.066) and 48\% for Llama (0.446 $\rightarrow$ 0.230). The consistent direction and magnitude across both models supports Retrieval Noise as a genuine, retrieval-driven failure mode rather than a model-specific artifact.

\textbf{Evidence Ignorance.} Moving the gold passage from top position to the middle of a 10-document context leaves token F1 essentially unchanged for both models (Claude: 0.067 $\rightarrow$ 0.067; Llama: 0.189 $\rightarrow$ 0.182), while removing it entirely causes a much larger drop (Claude: 0.067 $\rightarrow$ 0.032; Llama: 0.182 $\rightarrow$ 0.063) that closely tracks the drop in retrieval recall (1.000 $\rightarrow$ 0.000). This pattern is more consistent with an \textit{evidence-availability} effect (the gold fact is present or absent) than with a \textit{lost-in-the-middle position} effect for this dataset and these two models at this context length (10 documents): burying the correct passage in the middle of the context did not measurably hurt either model, but deleting it did. Our manual error analysis in Section~\ref{sec:empirical-erroranalysis} identifies a small number of individual questions where position \textit{did} matter (correct at top/middle, wrong when removed), alongside others where the model's answer was identical regardless of position, suggesting the aggregate pattern combines both a real evidence-availability effect and item-specific reading-comprehension behavior that is not purely positional.

\textbf{Trial-to-trial variance.} For the four three-trial stress conditions, per-trial mean token F1 varies by at most $\pm 0.001$--$0.004$ for Claude and $\pm 0.003$--$0.012$ for Llama (population standard deviation across the three trial-level means; full per-condition figures available in the released logs). This is small relative to the clean-vs-stress effect sizes reported above, consistent with temperature 0.0 producing near-deterministic outputs for both models, with a small residual amount of run-to-run variation likely attributable to backend batching or floating-point non-associativity on the provider side rather than sampling randomness.

\subsection{Manual Error Analysis}
\label{sec:empirical-erroranalysis}

To go beyond aggregate metrics, we drew a stratified sample of 70 generation outputs (5 questions $\times$ 7 mode/condition combinations $\times$ 2 models) and manually categorized each into one of: \textit{grounded\_correct}, \textit{grounded\_incorrect}, \textit{ungrounded\_correct}, \textit{partial\_grounding}, \textit{refused\_abstained}, or \textit{retrieval\_failure\_not\_model\_fault}, following a fixed labeling rubric. \textbf{This labeling pass was first drafted by an AI assistant (this paper's Computer-based research assistant), then independently audited by the human author against a blind copy of the same 70 items (no access to the AI-drafted labels during labeling).} A human audit of this full subsample found substantial agreement with the AI-assisted first-pass labels (Cohen's $\kappa = 0.71$, 81\% exact agreement across all 70 items); we report the raw labeled sample (questions, answers, retrieved passage IDs, both label sets, and rater notes) alongside this paper for independent verification.

Across the 70-item sample, 39 items (56\%) were \textit{grounded\_correct}, 15 (21\%) \textit{grounded\_incorrect}, 7 (10\%) \textit{refused\_abstained}, 4 (6\%) \textit{retrieval\_failure\_not\_model\_fault}, 3 (4\%) \textit{ungrounded\_correct}, and 2 (3\%) \textit{partial\_grounding} (AI-assisted first-pass distribution). Per model, Claude's 35 sampled items split 20 \textit{grounded\_correct}, 8 \textit{grounded\_incorrect}, 3 \textit{ungrounded\_correct}, 2 \textit{retrieval\_failure\_not\_model\_fault}, 2 \textit{refused\_abstained}; Llama's 35 split 19 \textit{grounded\_correct}, 7 \textit{grounded\_incorrect}, 5 \textit{refused\_abstained}, 2 \textit{retrieval\_failure\_not\_model\_fault}, 2 \textit{partial\_grounding}. The independent human audit shifted 13 of the 70 items to a different label, mostly toward \textit{partial\_grounding} for ambiguous cases the AI pass had split into \textit{retrieval\_failure\_not\_model\_fault} or \textit{refused\_abstained} (full human-corrected distribution released alongside this paper); the three qualitative patterns below were individually re-checked against the human-corrected labels and hold under both label sets. This sample is far too small to support a statistical claim that one model is generally ``better calibrated'' than the other, but three qualitative patterns recurred often enough across matched items (same question, same condition, both models) to be worth reporting as hypotheses for future, larger-scale study:

\begin{itemize}[leftmargin=*,nosep]
\item \textbf{Retrieval-driven failures replicate across models.} On several Retrieval Noise items where the gold passage was displaced from the top-5 by distractors, both Claude and Llama failed or partially failed on the identical item, and both succeeded on the matched clean-condition counterpart. This is direct evidence that these specific failures are attributable to the retrieval stage, not to either model's language understanding.
\item \textbf{Calibration under missing evidence differed between models on some items.} On two matched Retrieval-Blindness/Evidence-Ignorance-removed items, Llama abstained (``I don't know based on the available information'') where Claude committed to an answer despite explicitly noting the context did not directly support it---in one case Claude's guess happened to be correct (an \textit{ungrounded\_correct} case), in another it was confidently wrong. This suggests Claude was, in this small sample, more willing to supplement missing retrieved evidence with an inference or parametric guess, while Llama more often deferred to the abstention behavior requested by the prompt template.
\item \textbf{Some failures were retrieval-quality artifacts independent of the stress condition.} For at least one item, the top-5 retrieved passages were identical between the clean and blind conditions, indicating imperfect retrieval recall (not the deliberate evidence-withholding manipulation) was the actual limiting factor for that item in both conditions.
\end{itemize}

We emphasize these are illustrative observations from 70 items, not validated findings; even with a human-audited label set, the sample size is too small to generalize beyond these specific items.

\subsection{Limitations of This Empirical Validation}
\label{sec:empirical-limitations}

\begin{itemize}[leftmargin=*,nosep]
\item \textbf{Two of three planned models.} GPT-4.1 was part of the original three-model design but is not included in this version. We hit a sequence of access obstacles unrelated to the taxonomy itself (billing-flow issues on one host, quota and model-availability changes on others). Given the time constraints of this study, we proceeded with the two available models and documented GPT-4.1 as explicit future work, rather than delay the study indefinitely or substitute an unplanned third model.
\item \textbf{Two failure modes not empirically tested.} Evidence Over-Reliance and Latency-Induced Failure remain conceptual claims in this paper, illustrated only by the enterprise scenarios in Section~\ref{sec:casestudy}. We did not design stress conditions or collect real data for either mode.
\item \textbf{Small sample size.} 30 base questions per failure mode, drawn from two public dataset families, is a preliminary sample, not a benchmark-scale evaluation. Effect sizes reported here should be read as directional evidence that the failure modes are real and measurable, not as precise, generalizable estimates.
\item \textbf{Single dataset family per mode.} RGB and Lost-in-the-Middle are widely used, well-documented public benchmarks, but each contributes a specific style of question and passage structure. We did not test across multiple corpora or domains within a mode.
\item \textbf{Exact match undercounts correctness on this dataset combination}, as discussed in Section~\ref{sec:empirical-design}; token F1 is a more reliable relative-comparison metric here, but neither is a substitute for semantic or human judgment of correctness.
\item \textbf{Manual error analysis covers only 70 of 900 generation calls.} The AI-assisted first-pass labels for this subsample have since been independently audited by the human author ($\kappa = 0.71$, 81\% exact agreement; see Section~\ref{sec:empirical-erroranalysis}), but the 70-item sample itself remains small. We still recommend a larger-scale, human-labeled error analysis covering more of the 900 generation calls before drawing definitive conclusions from the qualitative patterns identified in Section~\ref{sec:empirical-erroranalysis}.
\item \textbf{Infrastructure history.} The Llama runs were migrated across three API hosts during data collection (an initially free-tier host that later removed the model from its catalog entirely; a second host blocked on account funding; a third, OpenRouter, that ultimately hosted the full run). All 450 Llama calls reported in Table~\ref{tab:llama-results} are from the final, single consistent host and model identifier; no results from the abandoned hosts are included as evidence.
\item \textbf{Single-provider access.} Both models were accessed through OpenRouter's API pass-through rather than each vendor's native API. We have no evidence this affects model outputs, and OpenRouter's documented pricing for Claude Sonnet 4.6 matches the native Anthropic list price with no added per-token markup, but we note the access path for reproducibility.
\end{itemize}

These limitations mean the results in this section should be read as a preliminary, defensible demonstration that three of our five proposed failure modes are real, measurable, and reproducible for two specific models under controlled conditions---not as a complete empirical validation of the taxonomy.

% ========== SECTION 5: CASE STUDY ==========
\section{Case Study: Enterprise RAG Failures}
\label{sec:casestudy}

To illustrate how our taxonomy and diagnosis guide apply in practice, we consider a simplified enterprise RAG system deployed for internal question answering over company documentation. The system indexes policy documents, technical manuals, and meeting notes, and exposes a chat interface to employees. These scenarios are illustrative rather than tied to a single deployed product, but they reflect patterns commonly reported in enterprise RAG pipelines.

\subsection{Scenario 1: Missing Risk Policy}

An employee asks: ``What is our policy on third-party security audits?'' The system responds with a generic description of security best practices but cites only high-level company values pages.

\textbf{Observation.} Retrieval logs show that the index contains a detailed ``Third-Party Risk Policy'' document, but it was not retrieved. The answer is vague and lacks specific policy clauses.

\textbf{Diagnosis.} The symptom maps to \textit{Retrieval Blindness}: the retriever failed to surface the relevant policy despite its presence in the corpus.

\textbf{Mitigation.} Engineers expand the corpus metadata, fine-tune embeddings on internal policy text, and add query rewriting templates for compliance-related questions. After these changes, the same query retrieves and cites the correct policy document.

\subsection{Scenario 2: Mixed Context on Benefits}

A query asks: ``How many vacation days do full-time employees get after three years?'' The system retrieves both current and outdated HR documents, as well as informal meeting notes. The answer correctly states the current number of days but also mentions an old policy and informal exceptions.

\textbf{Observation.} The retrieved context includes off-topic fragments and outdated information that the model blends into a single answer.

\textbf{Diagnosis.} This pattern matches \textit{Retrieval Noise} combined with \textit{Evidence Over-Reliance}: the retriever returns noisy and outdated content, and the model trusts it without filtering.

\textbf{Mitigation.} Engineers add stricter metadata filters (time stamps and document type), introduce a re-ranking step to prioritize official HR policies, and mark archived documents as low-priority. They also add a verification step that cross-checks answers against the most recent policy document.

\subsection{Scenario 3: Ignored Incident Report}

An incident response team asks: ``What were the main root causes of the outage last quarter?'' The system retrieves a detailed incident report but answers with a generic explanation about ``network instability'' that does not match the report's specific causes.

\textbf{Observation.} Removing the incident report from the context barely changes the answer, indicating that the model is not grounding its response in the retrieved evidence.

\textbf{Diagnosis.} This is a clear case of \textit{Evidence Ignorance}: retrieval succeeds, but the model ignores the relevant context, relying on parametric priors instead.

\textbf{Mitigation.} Engineers strengthen instructions to explicitly summarize and cite key passages from incident reports, reduce context length to avoid lost-in-the-middle effects, and place the most important sections at the beginning of the context window. Subsequent answers align more closely with the retrieved incident report.

\subsection{Scenario 4: Latency-Driven Shortcut Answers}

A regional operations team uses the RAG system during a live incident to ask: ``Which data centers are currently at risk if we shut down the west region for maintenance?'' Under normal load, the system retrieves capacity dashboards, failover playbooks, and recent incident reports before summarizing the risk. During peak traffic, however, the retrieval service is throttled and the orchestration layer enforces a strict latency budget. The system returns a short answer that mentions only the west region and offers generic reassurance, without referencing any of the actual capacity or failover documents.

\textbf{Observation.} Monitoring logs show that only the first hop of retrieval completed before the latency budget expired; downstream queries to capacity dashboards and incident reports were silently skipped. The generated answer is fast but incomplete, and it ignores critical context that would have changed the recommended action.

\textbf{Diagnosis.} This pattern matches \textit{Latency-Induced Failure}: retrieval and orchestration logic are constrained by time budgets and load, causing the system to answer from partial context and default assumptions, rather than fully grounded evidence.

\textbf{Mitigation.} Engineers profile the end-to-end latency budget and separate ``must-have'' retrieval steps (e.g., capacity dashboards and incident reports) from optional ones. They introduce caching for frequently accessed documents, relax timeouts for high-risk queries, and design a graceful degradation path that explicitly signals when an answer is based on partial context. They also add latency-aware logging so that incident retrospectives can distinguish between failures caused by retrieval quality and failures caused by time-constrained context.

These four scenarios demonstrate how a practitioner can use our taxonomy and diagnosis guide to label failures, understand their causes, and select targeted mitigations. While simplified, they reflect common patterns reported in production RAG deployments and show how a compact, five-mode taxonomy can structure failure analysis.

% ========== SECTION 6: IMPLICATIONS FOR EVALUATION ==========
\section{Implications for Evaluation}
\label{sec:evaluation}

Our taxonomy reveals a blind spot in current RAG evaluation. Benchmarks such as RAGAS and RGB primarily measure accuracy, faithfulness, and relevance under relatively clean retrieval conditions \cite{ragas2024,rgb2024}. A system may score highly while remaining vulnerable to the failures we identify.

We call for \textbf{failure-aware benchmarks} that intentionally inject each failure mode and report per-category resilience scores. Concretely, such benchmarks could:
\begin{itemize}[leftmargin=*,nosep]
\item Construct scenarios where the corpus is incomplete (Retrieval Blindness) or deliberately noisy (Retrieval Noise).
\item Evaluate whether models rely on or ignore context (Evidence Ignorance and Evidence Over-Reliance) under conflicting or adversarial documents.
\item Measure performance under constrained latency budgets (Latency-Induced Failure), varying retrieval time and system load.
\end{itemize}

For example, a failure-aware benchmark could report retrieval recall and corpus coverage for Retrieval Blindness; precision at $k$, nDCG, and distractor sensitivity for Retrieval Noise; context-ablation sensitivity and faithfulness for Evidence Ignorance; conflict-resolution accuracy and citation correctness for Evidence Over-Reliance; and grounded-answer quality, timeout rate, and end-to-end latency under fixed retrieval budgets for Latency-Induced Failure.

These protocols would complement existing RAG evaluation suites by stressing robustness across failure modes, rather than only measuring correctness when retrieval behaves ideally.

% ========== SECTION 7: CONCLUSION ==========
\section{Conclusion}
\label{sec:conclusion}

We presented a taxonomy of five failure modes in RAG systems: Retrieval Blindness, Retrieval Noise, Evidence Ignorance, Evidence Over-Reliance, and Latency-Induced Failure. We contributed definitions, common causes, a practical diagnosis guide, and a call for failure-aware evaluation protocols. In this version, we took a first empirical step beyond the conceptual framework: a preliminary, controlled study of three of the five modes---Retrieval Blindness, Retrieval Noise, and Evidence Ignorance---using two fixed generator models (Claude Sonnet 4.6 and Llama 3.3 70B Instruct), an identical retrieval pipeline and evaluation protocol, and 900 real, logged API calls (Section~\ref{sec:empirical}). Both models showed substantial, directionally consistent degradation from clean to stress conditions across all three modes, and retrieval-quality metrics were identical between models within each condition as expected, since retrieval is a property of the shared pipeline rather than the generator. This supports the taxonomy's core claims for these three modes, for these two models, on this dataset---while leaving Evidence Over-Reliance, Latency-Induced Failure, a third model (GPT-4.1), and broader dataset/domain coverage as open empirical questions.

By offering a shared vocabulary, a reusable diagnostic tool, and now a reproducible (if preliminary) empirical basis, we aim to make this taxonomy a standard reference for describing, debugging, and testing RAG failures. The four enterprise case study scenarios provide concrete illustrations of how the taxonomy and diagnosis guide apply in practice. Future work should prioritize: (1) completing the three-model design with GPT-4.1 once stable API access is available; (2) designing and running controlled experiments for Evidence Over-Reliance and Latency-Induced Failure, which remain conceptual in this paper; (3) scaling the empirical study to larger samples, additional datasets and domains, and a fully human-labeled error analysis beyond the 70-item AI-assisted-and-human-audited subsample reported here; and (4) building on the failure-aware evaluation protocols proposed in Section~\ref{sec:evaluation} to create dedicated benchmarks and open-source tooling that operationalize the diagnosis guide for real-world RAG pipelines. We release our experimental configuration, prompts, raw outputs, and error-analysis sample alongside this paper to support that future work.

% ========== BIBLIOGRAPHY ==========
\bibliographystyle{plain}
\bibliography{references}

\end{multicols}

\end{document}